\chapter{2007年四川卷（理）}
\section{单选题}
\begin{problem}
复数 \(\frac{1+\i}{1-\i} + \i^3\) 的值是（\quad）

\choices
    {\(0\)}
    {\(1\)}
    {\(-1\)}
    {\(1\)}
\end{problem}

\answer{A}

\begin{solution}
\[
\frac{1+\i}{1-\i}=\frac{(1+\i)^2}{(1-\i)(1+\i)}=\frac{2\i}{2}=\i
\]
又因为 \(\i^3=-\i\)，所以原式 \(=\i-\i=0\)，故选 A.
\end{solution}

\begin{problem}
函数 \( f(x) = 1 + \log_2 x \) 与 \( g(x) = 2^{-x + 1} \) 在同一直角坐标系下的图象大致是（\quad）

\choices
    {\choicebitmap[width=0.15\paperwidth]{img/2007/sichuan_science/q2_split_optA.png}}
    {\choicebitmap[width=0.15\paperwidth]{img/2007/sichuan_science/q2_split_optB.png}}
    {\choicebitmap[width=0.15\paperwidth]{img/2007/sichuan_science/q2_split_optC.png}}
    {\choicebitmap[width=0.15\paperwidth]{img/2007/sichuan_science/q2_split_optD.png}}
\end{problem}

\answer{C}

\begin{solution}
函数 \(f(x)=1+\log_2x\) 的定义域为 \(x>0\)，且底数 \(2>1\)，所以其图象在定义域内单调递增，经过点 \((1,1)\)，并以 \(x=0\) 为竖直渐近线. 函数 \(g(x)=2^{1-x}\) 始终为正，在整个定义域上单调递减，经过点 \((0,2)\)，\((1,1)\)，并以 \(y=0\) 为水平渐近线. 两函数一个递增，一个递减，且都经过 \((1,1)\)，因此图象只能是选项 C，故选 C.
\end{solution}

\begin{problem}
\(\displaystyle \lim\limits_{x\to 1}\frac{x^2 - 1}{2x^2 - x - 1} =\)（\quad）

\choices
    {\(0\)}
    {\(1\)}
    {\( \frac{1}{2} \)}
    {\( \frac{2}{3} \)}
\end{problem}

\answer{D}

\begin{solution}
\[
2x^2-x-1=(2x+1)(x-1)
\]
\[
x^2-1=(x-1)(x+1)
\]
当 \(x\ne1\) 时，
\[
\frac{x^2-1}{2x^2-x-1}=\frac{x+1}{2x+1}
\]
所以
\[
\lim_{x\to1}\frac{x^2-1}{2x^2-x-1}=\frac{1+1}{2+1}=\frac23
\]
故选 D.
\end{solution}

\begin{problem}
如图，\(ABCD - A_{1}B_{1}C_{1}D_{1}\) 为正方体，下面结论错误的是（\quad）

\bitmapfigure[max width=0.34\linewidth]{img/2007/sichuan_science/q4_fig1.png}

\choices
    {\(BD \parallel\) 平面 \(CB_{1}D_{1}\)}
    {\(AC_{1} \perp BD\)}
    {\(AC_{1} \perp\) 平面 \(CB_{1}D_{1}\)}
    {异面直线 \(AD\) 与 \(CB_{1}\) 角为 \(60^{\circ}\)}
\end{problem}

\answer{D}

\begin{solution}
取正方体棱长为 \(1\)，建立直角坐标系，令 \(A(0,0,0)\)，\(B(1,0,0)\)，\(C(1,1,0)\)，\(D(0,1,0)\)，\(A_1(0,0,1)\)，\(B_1(1,0,1)\)，\(C_1(1,1,1)\)，\(D_1(0,1,1)\).
平面 \(CB_1D_1\) 的方程为 \(x+y+z=2\)，其法向量可取 \(\bs n=(1,1,1)\).

有 \(\overrightarrow{BD}=(-1,1,0)\)，所以 \(\overrightarrow{BD}\cdot\bs n=0\)，且直线 \(BD\) 上各点满足 \(x+y+z=1\)，故 \(BD\parallel\) 平面 \(CB_1D_1\)，选项 A 正确.

又 \(\overrightarrow{AC_1}=(1,1,1)\)，\(\overrightarrow{BD}=(-1,1,0)\)，从而
\[
\overrightarrow{AC_1}\cdot\overrightarrow{BD}=0
\]
选项 B 正确；因为 \(\overrightarrow{AC_1}\) 与平面 \(CB_1D_1\) 的法向量平行，所以 \(AC_1\perp\) 平面 \(CB_1D_1\)，选项 C 正确.

最后，\(\overrightarrow{AD}=(0,1,0)\)，\(\overrightarrow{CB_1}=(0,-1,1)\)，两异面直线的夹角取锐角，故
\[
\cos\theta=\frac{|\overrightarrow{AD}\cdot\overrightarrow{CB_1}|}{|\overrightarrow{AD}|\,|\overrightarrow{CB_1}|}=\frac1{\sqrt2}
\]
即 \(\theta=45^\circ\)，不是 \(60^\circ\). 故错误的是 D.
\end{solution}

\begin{problem}
如果双曲线 \(\frac{x^2}{4} - \frac{y^2}{2} = 1\) 上一点 \(P\) 到双曲线右焦点的距离是 \(2\)，那么点 \(P\) 到 \(y\) 轴的距离是（\quad）

\choices
    {\(\frac{4\sqrt{6}}{3}\)}
    {\(\frac{2\sqrt{6}}{3}\)}
    {\(2\sqrt{6}\)}
    {\(2\sqrt{3}\)}
\end{problem}

\answer{A}

\begin{solution}
双曲线可写为 \(x^2-2y^2=4\)，其右焦点为 \(F(\sqrt6,0)\). 设 \(P(x,y)\)，由 \(PF=2\) 得
\[
(x-\sqrt6)^2+y^2=4
\]
由双曲线方程得 \(y^2=\frac{x^2-4}{2}\)，代入上式：
\[
x^2-2\sqrt6x+6+\frac{x^2-4}{2}=4
\]
即
\[
\frac32x^2-2\sqrt6x=x\left(\frac32x-2\sqrt6\right)=0
\]
\(x=0\) 不满足双曲线方程，故 \(x=\frac{4\sqrt6}{3}>0\). 点 \(P\) 到 \(y\) 轴的距离为 \(|x|=\frac{4\sqrt6}{3}\)，故选 A.
\end{solution}

\begin{problem}
设球 \(O\) 的半径是 \(1\)，\(A\)，\(B\)，\(C\) 是球面上三点，已知 \(A\) 到 \(B\)，\(C\) 两点的球面距离都是 \(\frac{\pi}{2}\)，且二面角 \(B - OA - C\) 的大小为 \(\frac{\pi}{3}\)，则从 \(A\) 点沿球面经 \(B\)，\(C\) 两点再回到 \(A\) 点的最短距离是（\quad）

\bitmapfigure[max width=0.30\linewidth]{img/2007/sichuan_science/q6_fig1.png}

\choices
    {\(\frac{7\pi}{6}\)}
    {\(\frac{5\pi}{4}\)}
    {\(\frac{4\pi}{3}\)}
    {\(\frac{3\pi}{2}\)}
\end{problem}

\answer{C}

\begin{solution}
球半径为 \(1\)，且 \(A\) 到 \(B\)，\(C\) 的球面距离均为 \(\frac\pi2\)，所以大圆上的弧 \(AB\)，\(AC\) 所对的圆心角均为 \(\frac\pi2\)，即 \(OA\perp OB\)，\(OA\perp OC\). 二面角 \(B-OA-C\) 是两个平面 \(BOA\)，\(COA\) 的夹角，在垂直于 \(OA\) 的截面内，它等于 \(OB\) 与 \(OC\) 的夹角，因此
\[
\angle BOC=\frac\pi3
\]
单位球面上两点间最短球面距离等于相应圆心角，所以从 \(A\) 经 \(B\)，\(C\) 回到 \(A\) 的最短距离为
\[
\frac\pi2+\frac\pi3+\frac\pi2=\frac{4\pi}{3}
\]
故选 C.
\end{solution}

\begin{problem}
设 \(A(a,1)\)，\(B(2,b)\)，\(C(4,5)\) 为坐标平面上三点，\(O\) 为坐标原点. 若 \(\overrightarrow{OA}\) 与 \(\overrightarrow{OB}\) 在 \(\overrightarrow{OC}\) 上的投影相同，则 \(a\) 与 \(b\) 满足的关系式为（\quad）

\choices
    {\(4a - 5b = 3\)}
    {\(5a - 4b = 3\)}
    {\(4a + 5b = 14\)}
    {\(5a + 4b = 14\)}
\end{problem}

\answer{A}

\begin{solution}
\(\overrightarrow{OA}\) 在 \(\overrightarrow{OC}\) 上的投影长度为 \(\dfrac{\overrightarrow{OA}\cdot\overrightarrow{OC}}{|\overrightarrow{OC}|}\)，\(\overrightarrow{OB}\) 的投影长度同理. 由于分母相同，投影相同等价于
\[
\overrightarrow{OA}\cdot\overrightarrow{OC}=\overrightarrow{OB}\cdot\overrightarrow{OC}
\]
代入 \(\overrightarrow{OA}=(a,1)\)，\(\overrightarrow{OB}=(2,b)\)，\(\overrightarrow{OC}=(4,5)\)，得
\[
4a+5=8+5b
\]
即 \(4a-5b=3\)，故选 A.
\end{solution}

\begin{problem}
已知抛物线 \( y = -x^2 + 3 \) 上存在关于直线 \( x + y = 0 \) 对称的相异两点 \(A\)，\(B\)，则 \(|AB|\) 等于（\quad）

\choices
    {\(3\)}
    {\(4\)}
    {\(3\sqrt{2}\)}
    {\(4\sqrt{2}\)}
\end{problem}

\answer{C}

\begin{solution}
设 \(A=(x,y)\) 在抛物线 \(y=-x^2+3\) 上. 关于直线 \(x+y=0\) 的对称变换为
\[
(x,y)\longmapsto(-y,-x)
\]
故其对称点 \(B=(-y,-x)\) 也在该抛物线上，由 \(A\)，\(B\) 均在抛物线上得 \(y=-x^2+3\)，\(-x=-y^2+3\)，即 \(x=y^2-3\). 两式相加得
\[
x+y=y^2-x^2=(y-x)(x+y)
\]
所以 \((x+y)(1-y+x)=0\). 若 \(x+y=0\)，则点在对称轴上，会与其对称点重合，与 \(A\)，\(B\) 相异矛盾. 因此 \(y-x=1\).
代入 \(y=-x^2+3\)，得 \(x+1=-x^2+3\)，即 \(x^2+x-2=0\).
解得 \(x=1\) 或 \(x=-2\)，对应两点为 \((1,2)\)，\((-2,-1)\). 于是
\[
|AB|=\sqrt{(1+2)^2+(2+1)^2}=3\sqrt2
\]
故选 C.
\end{solution}

\begin{problem}
某公司有 \(60\text{ 万元}\) 资金，计划投资甲，乙两个项目，按要求对项目甲的投资不小于对项目乙投资的 \(\frac{2}{3}\) 倍，且对每个项目的投资不能低于 \(5\text{ 万元}\)，对项目甲每投资 \(1\text{ 万元}\) 可获得 \(0.4\text{ 万元}\) 的利润，对项目乙每投资 \(1\text{ 万元}\) 可获得 \(0.6\text{ 万元}\) 的利润，该公司正确规划投资后，在这两个项目上共可获得的最大利润为（\quad）

\choices
    {\(36\text{ 万元}\)}
    {\(31.2\text{ 万元}\)}
    {\(30.4\text{ 万元}\)}
    {\(24\text{ 万元}\)}
\end{problem}

\answer{B}

\begin{solution}
设投资甲，乙项目分别为 \(x\)，\(y\) 万元，则约束条件为 \(x+y\leq60\)，\(x\geq\frac23y\)，\(x\geq5\)，\(y\geq5\).
总利润为
\[
P=0.4x+0.6y=0.4(x+y)+0.2y
\]
由 \(x\geq\frac23y\) 得 \(x+y\geq\frac53y\)，再结合 \(x+y\leq60\)，可得 \(y\leq36\). 因此
\[
P\leq0.4\times60+0.2\times36=31.2
\]
当 \(x=24\)，\(y=36\) 时，所有约束均满足且上式取等号，所以最大利润为 \(31.2\) 万元，故选 B.
\end{solution}

\begin{problem}
用数字 \(0\)，\(1\)，\(2\)，\(3\)，\(4\)，\(5\) 可以组成没有重复数字，并且比 \(20000\) 大的五位偶数共有（\quad）

\choices
    {\(288\text{ 个}\)}
    {\(240\text{ 个}\)}
    {\(144\text{ 个}\)}
    {\(126\text{ 个}\)}
\end{problem}

\answer{B}

\begin{solution}
五位数的首位只能取 \(2\)，\(3\)，\(4\)，\(5\)，末位必须取偶数字 \(0\)，\(2\)，\(4\)，且各位数字不能重复.

首位为 \(2\) 时，末位不能为 \(2\)，末位取 \(0\) 或 \(4\)，每种情形中间三位从余下四个数字中选排，均有 \(\mathrm A_4^3=24\) 个，共 \(2\times24=48\) 个. 首位为 \(3\) 时末位有 \(0\)，\(2\)，\(4\) 三种选择，共 \(3\times24=72\) 个. 首位为 \(4\) 时末位有 \(0\)，\(2\) 两种选择，共 \(2\times24=48\) 个. 首位为 \(5\) 时末位有三种选择，共 \(3\times24=72\) 个.

所以共有
\[
48+72+48+72=240
\]
个，故选 B.
\end{solution}

\begin{problem}
如图，\(l_{1}\)，\(l_{2}\)，\(l_{3}\) 是同一平面内的三条平行直线，\(l_{1}\) 与 \(l_{2}\) 间的距离是 \(1\)，\(l_{2}\) 与 \(l_{3}\) 间的距离是 \(2\)，正三角形 \(ABC\) 的三顶点分别在 \(l_{1}\)，\(l_{2}\)，\(l_{3}\) 上，则 \(\triangle ABC\) 的边长是（\quad）

\bitmapfigure[max width=0.36\linewidth]{img/2007/sichuan_science/q11_fig1.png}

\choices
    {\(2\sqrt{3}\)}
    {\(\frac{4\sqrt{6}}{3}\)}
    {\(\frac{3-\sqrt{7}}{4}\)}
    {\(\frac{2\sqrt{21}}{3}\)}
\end{problem}

\answer{D}

\begin{solution}
设 \(l_2\) 为 \(x\) 轴，取垂直于三条平行线的方向为纵轴，使得 \(l_1\) 的方程为 \(y=1\)，\(l_2\) 的方程为 \(y=0\)，\(l_3\) 的方程为 \(y=-2\).
设 \(B=(0,0)\)，\(A=(p,1)\)，\(C=(q,-2)\)，并记正三角形边长为 \(s\). 由 \(AB=BC=s\)，有 \(p^2+1=s^2\)，\(q^2+4=s^2\)，从而 \(p^2-q^2=3\).
在等边三角形中，\(\angle ABC=60^\circ\)，所以
\[
(p,1)\cdot(q,-2)=\frac{s^2}{2}
\]
即 \(pq-2=\frac{s^2}{2}\). 使用 \(s^2=p^2+1\)，得
\[
2pq=p^2+5
\]
令 \(t=p^2\)，由 \(q^2=p^2-3=t-3\)，对上式平方得
\[
4t(t-3)=(t+5)^2
\]
即 \(3t^2-22t-25=0\). 解得 \(t=\frac{25}{3}\) 或 \(t=-1\)，舍去负值.
所以 \(s^2=p^2+1=\frac{28}{3}\)，从而 \(s=\frac{2\sqrt{21}}{3}\).
故选 D.
\end{solution}

\begin{problem}
已知一组抛物线 \( y = \frac{1}{2}ax^2 + bx + 1 \)，其中 \(a\) 为 \(2\)，\(4\)，\(6\)，\(8\) 中任取的一个数，\(b\) 为 \(1\)，\(3\)，\(5\)，\(7\) 中任取的一个数. 从这些抛物线中任意抽取两条，它们在与直线 \(x = 1\) 交点处的切线相互平行的概率是（\quad）

\choices
    {\(\frac{1}{12}\)}
    {\(\frac{7}{60}\)}
    {\(\frac{6}{25}\)}
    {\(\frac{5}{25}\)}
\end{problem}

\answer{B}

\begin{solution}
每条抛物线在 \(x=1\) 处的切线斜率为 \(y'=ax+b\)，所以 \(y'(1)=a+b\).
共有 \(4\times4=16\) 条抛物线，任取两条的总数为 \(\mathrm C_{16}^{2}=120\). 由于 \(a\) 取 \(2\)，\(4\)，\(6\)，\(8\)，\(b\) 取 \(1\)，\(3\)，\(5\)，\(7\)，和 \(a+b\) 依次为 \(3\)，\(5\)，\(7\)，\(9\)，\(11\)，\(13\)，\(15\) 时，对应的抛物线条数分别为 \(1\)，\(2\)，\(3\)，\(4\)，\(3\)，\(2\)，\(1\).

因此斜率相同的两条抛物线共有
\[
0+1+3+6+3+1+0=14
\]
对，所求概率为
\[
\frac{14}{\mathrm C_{16}^{2}}=\frac{14}{120}=\frac7{60}
\]
故选 B.
\end{solution}

\section{填空题}
\begin{problem}
设函数 \(f(x) = \e^{-(x-u)^2}\)（\(\e\) 是自然对数的底数）的最大值是 \(m\)，且 \(f(x)\) 是偶函数，则 \(m + u =\fillinblank{}\).
\end{problem}

\answer{\(1\)}

\begin{solution}
因为 \((x-u)^2\geq0\)，所以
\[
f(x)=\e^{-(x-u)^2}\leq1
\]
当 \(x=u\) 时取等号，故最大值 \(m=1\). 又因为 \(f(x)\) 为偶函数，特别有 \(f(1)=f(-1)\)，从而
\[
\e^{-(1-u)^2}=\e^{-(-1-u)^2}
\]
即 \((1-u)^2=(1+u)^2\)，解得 \(u=0\). 因此 \(m+u=1\).
\end{solution}

\begin{problem}
如图，在正三棱柱 \(ABC - A_{1}B_{1}C_{1}\) 中，侧棱长为 \(\sqrt{2}\)，底面三角形的边长为 \(1\)，则 \(BC_{1}\) 与侧面 \(ACC_{1}A_{1}\) 所成的角是\fillinblank{}.

\bitmapfigure[max width=6.0cm]{img/2007/sichuan_science/q14_fig1.png}
\end{problem}

\answer{\(\frac{\pi}{6}\)}

\begin{solution}
建立直角坐标系，令 \(A=(0,0,0)\)，\(B=(1,0,0)\)，\(C=\left(\frac12,\frac{\sqrt3}{2},0\right)\).
并令棱柱的高为 \(\sqrt2\)，则
\[
C_1=\left(\frac12,\frac{\sqrt3}{2},\sqrt2\right)
\]
平面 \(ACC_1A_1\) 由向量 \(\overrightarrow{AC}=\left(\frac12,\frac{\sqrt3}{2},0\right)\) 和竖直向量张成，其一个单位法向量为
\[
\bs n=\left(\frac{\sqrt3}{2},-\frac12,0\right)
\]
直线 \(BC_1\) 的方向向量为
\[
\bs v=\overrightarrow{BC_1}=\left(-\frac12,\frac{\sqrt3}{2},\sqrt2\right)
\]
且 \(|\bs v|=\sqrt3\). 设直线与平面所成的锐角为 \(\theta\)，则
\[
\sin\theta=\frac{|\bs v\cdot\bs n|}{|\bs v|\,|\bs n|}
=\frac{\left|-\frac{\sqrt3}{4}-\frac{\sqrt3}{4}\right|}{\sqrt3}=\frac12
\]
所以 \(\theta=\frac\pi6\).
\end{solution}

\begin{problem}
已知 \(\odot O\) 的方程是 \(x^2 + y^2 - 2 = 0\)，\(\odot O'\) 的方程是 \(x^2 + y^2 - 8x + 10 = 0\)，由动点 \(P\) 向 \(\odot O\) 和 \(\odot O'\) 所引的切线长相等，则动点 \(P\) 的轨迹方程是\fillinblank{}.
\end{problem}

\answer{\(x = \frac{3}{2}\)}

\begin{solution}
将两圆方程化为标准形式：
\[
x^2+y^2=2
\]
和
\[
(x-4)^2+y^2=6
\]
设动点 \(P=(x,y)\)，则从 \(P\) 向两圆所引切线长的平方分别为
\[
PT^2=x^2+y^2-2
\]
以及
\[
PT'^2=(x-4)^2+y^2-6=x^2+y^2-8x+10
\]
切线长相等，故
\[
x^2+y^2-2=x^2+y^2-8x+10
\]
解得 \(x=\frac32\). 当 \(x=\frac32\) 时两切线长平方均为 \(y^2+\frac14>0\)，所以轨迹为直线
\[
x=\frac32
\]
\end{solution}

\begin{problem}
下面有五个命题：

\begin{circlelist}
\item 函数 \(y = \sin^{4}x - \cos^{4}x\) 的最小正周期是 \(\pi\)；
\item 终边在 \(y\) 轴上的角的集合是 \(\{\alpha \mid \alpha = \frac{k\pi}{2}, k\in\Z\}\)；
\item 在同一坐标系中，函数 \(y = \sin x\) 的图象和函数 \(y = x\) 的图象有三个公共点；
\item 把函数 \(y = 3\sin\!\left(2x + \frac{\pi}{3}\right)\) 的图象向右平移 \(\frac{\pi}{6}\) 得到 \(y = 3\sin 2x\) 的图象；
\item 函数 \(y = \sin\!\left(x - \frac{\pi}{2}\right)\) 在 \([0,\pi]\) 上是减函数.
\end{circlelist}

其中真命题的序号是\fillinblank{}.（写出所有）
\end{problem}

\answer{\(\circled{1}\circled{4}\)}

\begin{solution}
命题 \(\circled{1}\) 中，
\[
\sin^4x-\cos^4x=(\sin^2x-\cos^2x)(\sin^2x+\cos^2x)=-\cos2x
\]
其最小正周期为 \(\pi\)，所以命题 \(\circled{1}\) 为真.

命题 \(\circled{2}\) 中，终边在 \(y\) 轴上的角应满足 \(\alpha=\frac\pi2+k\pi\)，而题中集合还包含 \(\alpha=0\) 等终边在 \(x\) 轴上的角，所以命题 \(\circled{2}\) 为假.

命题 \(\circled{3}\) 中，公共点的横坐标满足 \(\sin x=x\). 由于 \(|\sin x|\leq1\)，只需考察 \([-1,1]\). 令 \(h(x)=\sin x-x\)，则 \(h'(x)=\cos x-1\leq0\)，且 \(h(0)=0\)，故在该区间内只有 \(x=0\) 一个解，不是三个公共点，命题 \(\circled{3}\) 为假.

命题 \(\circled{4}\) 中，函数图象向右平移 \(\frac\pi6\) 后，横坐标代换为 \(x-\frac\pi6\)，得到
\[
3\sin\left(2\left(x-\frac\pi6\right)+\frac\pi3\right)=3\sin2x
\]
所以命题 \(\circled{4}\) 为真.

命题 \(\circled{5}\) 中，
\[
\sin\left(x-\frac\pi2\right)=-\cos x
\]
其导数为 \(\sin x\geq0\)（在 \((0,\pi)\) 内大于 \(0\)），故函数在 \([0,\pi]\) 上递增而非递减，命题 \(\circled{5}\) 为假. 因此真命题的序号为 \(\circled{1}\)，\(\circled{4}\).
\end{solution}

\section{解答题}
\begin{problem}
已知 \(\cos\alpha = \frac{1}{7}\)，\(\cos(\alpha - \beta) = \frac{13}{14}\)，且 \(0 < \beta < \alpha < \frac{\pi}{2}\).

\begin{enumerate}
\item 求 \(\tan 2\alpha\) 的值；
\item 求 \(\beta\).
\end{enumerate}
\end{problem}

\begin{solution}
\begin{enumerate}
\item 因为 \(0<\alpha<\frac{\pi}{2}\)，所以 \(\sin\alpha=\sqrt{1-\cos^2\alpha}=\frac{4\sqrt{3}}{7}\). 从而 \(\cos 2\alpha=2\cos^2\alpha-1=-\frac{47}{49}\)，\(\sin 2\alpha=2\sin\alpha\cos\alpha=\frac{8\sqrt{3}}{49}\)，故 \(\tan 2\alpha=\frac{\sin 2\alpha}{\cos 2\alpha}=-\frac{8\sqrt{3}}{47}\).
\item 令 \(\delta=\alpha-\beta\). 由 \(0<\beta<\alpha<\frac{\pi}{2}\) 可知 \(0<\delta<\frac{\pi}{2}\)，所以 \(\sin\delta=\sqrt{1-\cos^2\delta}=\frac{3\sqrt{3}}{14}\). 由差角公式，\(\cos\beta=\cos(\alpha-\delta)=\cos\alpha\cos\delta+\sin\alpha\sin\delta=\frac{1}{7}\cdot\frac{13}{14}+\frac{4\sqrt{3}}{7}\cdot\frac{3\sqrt{3}}{14}=\frac{1}{2}\). 又 \(0<\beta<\frac{\pi}{2}\)，故 \(\beta=\frac{\pi}{3}\).
\end{enumerate}
\end{solution}


\begin{problem}
厂家在产品出厂前，需对产品做检验，厂家将一批产品发给商家时，商家按合同规定也需随机抽取一定数量的产品做检验，以决定是否接收这批产品.

\begin{enumerate}
\item 若厂家库房中的每件产品合格的概率为 \(0.8\)，从中任意取出 \(4\) 件进行检验. 求至少有 \(1\) 件是合格品的概率；
\item 若厂家发给商家 \(20\) 件产品，其中有 \(3\) 件不合格，按合同规定该商家从中任取 \(2\) 件，都进行检验，只有 \(2\) 件都合格时才接收这批产品，否则拒收. 求该商家可能检验出不合格产品数 \(\xi\) 的分布列及期望 \(E(\xi)\)，并求该商家拒收这批产品的概率.
\end{enumerate}
\end{problem}

\begin{solution}
\begin{enumerate}
\item 设取出的 \(4\) 件产品中合格品数为 \(X\). 事件“至少有 \(1\) 件合格品”的对立事件是“\(4\) 件都不合格”，因此所求概率为 \(1-(1-0.8)^4=1-0.2^4=0.9984\).
\item 共有 \(\mathrm C_{20}^{2}=190\) 种等可能的取法，随机变量 \(\xi\) 的可能取值为 \(0\)，\(1\)，\(2\). 其中 \(3\) 件不合格，\(17\) 件合格，所以 \(P(\xi=0)=\frac{\mathrm C_{17}^{2}}{\mathrm C_{20}^{2}}=\frac{136}{190}\)，\(P(\xi=1)=\frac{\mathrm C_{3}^{1}\mathrm C_{17}^{1}}{\mathrm C_{20}^{2}}=\frac{51}{190}\)，\(P(\xi=2)=\frac{\mathrm C_{3}^{2}}{\mathrm C_{20}^{2}}=\frac{3}{190}\).
因此 \(E(\xi)=0\cdot\frac{136}{190}+1\cdot\frac{51}{190}+2\cdot\frac{3}{190}=\frac{3}{10}\). 只有当 \(\xi=0\) 时接收，所以拒收概率为 \(1-P(\xi=0)=1-\frac{136}{190}=\frac{27}{95}\).
\end{enumerate}
\end{solution}


\begin{problem}
如图，\(PCBM\) 是直角梯形，\(\angle PCB = 90^\circ\)，\(PM \parallel BC\)，\(PM = 1\)，\(BC = 2\)，又 \(AC = 1\)，\(\angle ACB = 120^\circ\)，\(AB \perp PC\)，直线 \(AM\) 与直线 \(PC\) 所成的角为 \(60^\circ\).

\begin{enumerate}
\item 求证：平面 \(PAC \perp\) 平面 \(ABC\)；
\item 求二面角 \(M - AC - B\) 的大小；
\item 求三棱锥 \(P - MAC\) 的体积.
\end{enumerate}

\bitmapfigure[max width=0.38\linewidth]{img/2007/sichuan_science/q19_fig1.png}
\end{problem}

\begin{solution}
\begin{enumerate}
\item 建立直角坐标系：以 \(C\) 为原点，令 \(CB\) 为 \(x\) 轴正向，\(PC\) 为 \(z\) 轴正向，并令 \(A\) 位于 \(y>0\) 的一侧. 由已知条件可设 \(B=(2,0,0)\)，\(P=(0,0,h)\)，\(M=(1,0,h)\). 因为 \(AB\perp PC\)，点 \(A\) 的 \(z\) 坐标为 \(0\)；又 \(AC=1\)，\(\angle ACB=120^{\circ}\)，故可取 \(A=\left(-\frac{1}{2},\frac{\sqrt{3}}{2},0\right)\). 于是 \(\overrightarrow{AM}=\left(\frac{3}{2},-\frac{\sqrt{3}}{2},h\right)\)，且 \(AM^2=3+h^2\). 由直线 \(AM\) 与直线 \(PC\) 所成的角为 \(60^{\circ}\)，得 \(\frac{h^2}{3+h^2}=\cos^2 60^{\circ}=\frac{1}{4}\)，所以 \(h^2=1\)，取 \(h=1\). 因为 \(PC\perp BC\)，且 \(PC\perp AB\)，而 \(AB\)，\(BC\) 是平面 \(ABC\) 内相交的两条直线，所以 \(PC\perp\) 平面 \(ABC\). 又 \(PC\subset\) 平面 \(PAC\)，故平面 \(PAC\perp\) 平面 \(ABC\).
\item 取平面 \(ABC\) 的法向量 \(\bs{n}_1=(0,0,1)\). 平面 \(MAC\) 内的两个向量可取 \(\overrightarrow{AC}=\left(\frac{1}{2},-\frac{\sqrt{3}}{2},0\right)\) 与 \(\overrightarrow{AM}=\left(\frac{3}{2},-\frac{\sqrt{3}}{2},1\right)\)，故其法向量可取
\[
\bs{n}_2=\overrightarrow{AC}\times\overrightarrow{AM}=\left(-\frac{\sqrt{3}}{2},-\frac{1}{2},\frac{\sqrt{3}}{2}\right).
\]
设二面角 \(M-AC-B\) 的大小为 \(\theta\)，则 \(0<\theta<\frac{\pi}{2}\)，从而 \(\cos\theta=\frac{|\bs{n}_1\cdot\bs{n}_2|}{|\bs{n}_1||\bs{n}_2|}=\sqrt{\frac{3}{7}}\). 因此 \(\tan\theta=\sqrt{\frac{1-3/7}{3/7}}=\frac{2\sqrt{3}}{3}\)，即 \(\theta=\arctan\frac{2\sqrt{3}}{3}\).
\item 由四面体体积的行列式公式，\(V_{P-MAC}=\frac{1}{6}\left|((M-C)\times(A-C))\cdot(P-C)\right|\). 代入 \(M-C=(1,0,1)\)，\(A-C=\left(-\frac{1}{2},\frac{\sqrt{3}}{2},0\right)\)，\(P-C=(0,0,1)\)，得 \(V_{P-MAC}=\frac{1}{6}\cdot\frac{\sqrt{3}}{2}=\frac{\sqrt{3}}{12}\).
\end{enumerate}
\end{solution}


\begin{problem}
设 \(F_{1}\)，\(F_{2}\) 分别是椭圆 \(\frac{x^{2}}{4} + y^{2} = 1\) 的左，右焦点.

\begin{enumerate}
\item 若 \(P\) 是该椭圆上的一个动点，求 \(\overrightarrow{PF_{1}} \cdot \overrightarrow{PF_{2}}\) 的最大值和最小值；
\item 设过定点 \(M(0,2)\) 的直线 \(l\) 与椭圆交于不同的两点 \(A\)，\(B\)，且 \(\angle AOB\) 为锐角（其中 \(O\) 为坐标原点），求直线 \(l\) 的斜率 \(k\) 的取值范围.
\end{enumerate}
\end{problem}

\begin{solution}
\begin{enumerate}
\item 椭圆方程可写为 \(x^2+4y^2=4\)，所以其半长轴为 \(2\)，半短轴为 \(1\)，焦半距为 \(\sqrt{2^2-1^2}=\sqrt{3}\). 因此 \(F_1=(-\sqrt{3},0)\)，\(F_2=(\sqrt{3},0)\). 设 \(P=(x,y)\)，则 \(\overrightarrow{PF_1}=(-\sqrt{3}-x,-y)\)，\(\overrightarrow{PF_2}=(\sqrt{3}-x,-y)\)，从而
\[
\overrightarrow{PF_1}\cdot\overrightarrow{PF_2}=(-\sqrt{3}-x)(\sqrt{3}-x)+y^2=x^2+y^2-3.
\]
由 \(x^2+4y^2=4\) 得 \(x^2=4-4y^2\)，所以 \(\overrightarrow{PF_1}\cdot\overrightarrow{PF_2}=1-3y^2\). 又 \(0\leqslant y^2\leqslant1\)，故其最大值为 \(1\)，最小值为 \(-2\).
\item 当直线 \(l\) 的斜率存在时，设其方程为 \(y=kx+2\). 将其代入椭圆方程，得 \((1+4k^2)x^2+16kx+12=0\). 设交点 \(A\)，\(B\) 的横坐标分别为 \(x_1\)，\(x_2\)，因为交于不同的两点，所以判别式满足
\[
\Delta=(16k)^2-4(1+4k^2)\cdot12=16(4k^2-3)>0,
\]
即 \(|k|>\frac{\sqrt{3}}{2}\). 由韦达定理，\(x_1+x_2=-\frac{16k}{1+4k^2}\)，\(x_1x_2=\frac{12}{1+4k^2}\). 又 \(A=(x_1,kx_1+2)\)，\(B=(x_2,kx_2+2)\)，所以
\[
\overrightarrow{OA}\cdot\overrightarrow{OB}=x_1x_2+(kx_1+2)(kx_2+2)=(1+k^2)x_1x_2+2k(x_1+x_2)+4=\frac{4(4-k^2)}{1+4k^2}.
\]
\(\angle AOB\) 为锐角等价于 \(\overrightarrow{OA}\cdot\overrightarrow{OB}>0\)，故 \(|k|<2\). 若 \(l\) 垂直于 \(x\) 轴，则 \(l\) 为 \(x=0\)，交点为 \((0,1)\)，\((0,-1)\)，此时 \(\angle AOB=\pi\)，不合题意. 综上，\(k\) 的取值范围为 \(\left(-2,-\frac{\sqrt{3}}{2}\right)\cup\left(\frac{\sqrt{3}}{2},2\right)\).
\end{enumerate}
\end{solution}

\begin{problem}
已知函数 \(f(x) = x^2 - 4\)，设曲线 \(y = f(x)\) 在点 \((x_n, f(x_n))\) 处的切线与 \(x\) 轴的交点为 \((x_{n+1}, 0)\)（\(n \in \N^{*}\)），其中 \(x_1\) 为正实数.

\begin{enumerate}
\item 用 \(x_n\) 表示 \(x_{n+1}\)；
\item 求证：对一切正整数 \(n\)，\(x_{n+1} \leqslant x_n\) 的充要条件是 \(x_1 \geqslant 2\)；
\item 若 \(x_1 = 4\)，记 \(a_n = \lg \frac{x_n + 2}{x_n - 2}\)，证明数列 \(\{a_n\}\) 成等比数列，并求数列 \(\{x_n\}\) 的通项公式.
\end{enumerate}
\end{problem}

\begin{solution}
\begin{enumerate}
\item 曲线在点 \((x_n,x_n^2-4)\) 处的切线斜率为 \(2x_n\)，其方程为 \(y-(x_n^2-4)=2x_n(x-x_n)\). 令 \(y=0\)，且 \(x=x_{n+1}\)，得 \(-(x_n^2-4)=2x_n(x_{n+1}-x_n)\)，所以 \(x_{n+1}=\frac{x_n^2+4}{2x_n}=\frac{1}{2}\left(x_n+\frac{4}{x_n}\right)\).
\item 由上式，\(x_{n+1}-x_n=\frac{4-x_n^2}{2x_n}\). 因为 \(x_1>0\)，且递推式保证各项均为正数. 若 \(x_1\geqslant2\)，则由均值不等式 \(x_{n+1}=\frac{1}{2}\left(x_n+\frac{4}{x_n}\right)\geqslant2\)，所以对每个正整数 \(n\) 都有 \(x_n\geqslant2\)，从而 \(x_{n+1}-x_n\leqslant0\)，即 \(x_{n+1}\leqslant x_n\). 反之，若对一切正整数 \(n\) 都有 \(x_{n+1}\leqslant x_n\)，取 \(n=1\)，由 \(x_2-x_1=\frac{4-x_1^2}{2x_1}\leqslant0\) 且 \(x_1>0\)，可得 \(x_1\geqslant2\). 因此，所述条件的充要条件是 \(x_1\geqslant2\).
\item 当 \(x_1=4\) 时，由第（2）问知 \(x_n\geqslant2\). 又由 \(x_1>2\) 及
\[
x_{n+1}-2=\frac{(x_n-2)^2}{2x_n}>0
\]
可知对一切正整数 \(n\) 都有 \(x_n>2\)，所以 \(a_n\) 有意义. 进一步，\(x_{n+1}+2=\frac{(x_n+2)^2}{2x_n}\)，从而
\[
\frac{x_{n+1}+2}{x_{n+1}-2}=\left(\frac{x_n+2}{x_n-2}\right)^2.
\]
两边取常用对数，得 \(a_{n+1}=2a_n\). 又 \(a_1=\lg\frac{4+2}{4-2}=\lg3\)，故 \(a_n=2^{n-1}\lg3\)，数列 \(\{a_n\}\) 是首项为 \(\lg3\)，公比为 \(2\) 的等比数列. 由 \(10^{a_n}=\frac{x_n+2}{x_n-2}\)，得 \(\frac{x_n+2}{x_n-2}=3^{2^{n-1}}\)，解得
\[
x_n=2\frac{3^{2^{n-1}}+1}{3^{2^{n-1}}-1}.
\]
\end{enumerate}
\end{solution}


\begin{problem}
设函数 \(f(x) = \left(1 + \frac{1}{n}\right)^x\)（\(n \in \N\)，且 \(n > 1\)，\(x \in \R\)）.

\begin{enumerate}
\item 当 \(x = 6\) 时，求 \(\left(1 + \frac{1}{n}\right)^x\) 的展开式中二项式系数最大的项；
\item 对任意的实数 \(x\)，证明 \(\frac{f(2x) + f(2)}{2} > f'(x)\)（\(f'(x)\) 是 \(f(x)\) 的导函数）；
\item 是否存在 \(a \in \N\)，使得 \(an < \sum_{k=1}^{n} \left(1 + \frac{1}{k}\right)^k < (a + 1)n\) 恒成立？若存在，试证明你的结论并求出 \(a\) 的值；若不存在，请说明理由.
\end{enumerate}
\end{problem}

\begin{solution}
\begin{enumerate}
\item 当 \(x=6\) 时，由二项式定理，\(\left(1+\frac{1}{n}\right)^6=1+\frac{6}{n}+\frac{15}{n^2}+\frac{20}{n^3}+\frac{15}{n^4}+\frac{6}{n^5}+\frac{1}{n^6}\). 其二项式系数依次为 \(1\)，\(6\)，\(15\)，\(20\)，\(15\)，\(6\)，\(1\)，最大值 \(20\) 只出现在第 \(4\) 项，所以二项式系数最大的项为 \(\frac{20}{n^3}\).
\item 令 \(r=1+\frac{1}{n}\)，则 \(r>1\)，\(f(x)=r^x\)，\(f(2x)=r^{2x}\)，\(f(2)=r^2\)，且 \(f'(x)=(\ln r)r^x\). 令 \(t=r^x>0\)，则由均值不等式，\(\frac{f(2x)+f(2)}{2}=\frac{t^2+r^2}{2}\geqslant rt\). 另一方面，\(\ln r=\int_{1}^{r}\frac{1}{s}\,\mathrm ds<\int_{1}^{r}1\,\mathrm ds=r-1<r\)，所以 \(rt>(\ln r)t=f'(x)\). 因而对任意实数 \(x\)，都有 \(\frac{f(2x)+f(2)}{2}>f'(x)\).
\item 题设中的 \(n\in\N\) 且 \(n>1\)，所以这里的 \(n\) 为不小于 \(2\) 的正整数. 记
\[
u_k=\left(1+\frac1k\right)^k
\]
当 \(k\geq2\) 时，由二项式定理
\[
u_k=1+1+\frac{\mathrm C_k^2}{k^2}+\cdots>2
\]
而 \(u_1=2\). 另一方面，对 \(j\geq2\)，有
\[
\frac{\mathrm C_k^j}{k^j}\leq\frac1{j!}\leq\frac1{2^{j-1}}
\]
所以
\[
u_k=2+\sum_{j=2}^{k}\frac{\mathrm C_k^j}{k^j}
\leq2+\sum_{j=2}^{k}\frac1{2^{j-1}}
=3-\frac1{2^{k-1}}<3
\]
于是对每个 \(n\geq2\)，
\[
2n=u_1+2(n-1)<\sum_{k=1}^{n}u_k<3n
\]
因此取 \(a=2\) 时题设不等式恒成立. 又取 \(n=2\) 时，
\[
\frac1n\sum_{k=1}^{n}u_k=\frac{2+\frac94}{2}=\frac{17}{8}
\]
若题设不等式对某个整数 \(a\) 恒成立，则 \(a<\frac{17}{8}<a+1\)，只能有 \(a=2\). 故存在这样的 \(a\)，且 \(a=2\).
\end{enumerate}
\end{solution}
